\documentclass[12pt, letterpaper]{article}

\usepackage[utf8]{inputenc}
\usepackage[T1]{fontenc}
\usepackage[margin=1in]{geometry}
\usepackage{setspace}
\usepackage{graphicx}
\usepackage{booktabs}
\usepackage{amsmath, amssymb}
\usepackage{hyperref}
\usepackage[style=apa, backend=biber]{biblatex}

\addbibresource{references.bib}
\onehalfspacing

\title{From Persons to Families to Organizations:\\A Cross-Level CAS Foundation for Brief Health Screening}
\author{Ryan Carr}
\date{\today}

\begin{document}

\maketitle

\begin{abstract}
TODO: Summarize the cross-level complex adaptive systems argument linking individual, family-system, and organizational health screening.
\end{abstract}

\section{Introduction}
TODO: Frame the need for a principled cross-level foundation rather than a single-domain analogy.

\section{Complex Adaptive Systems Across Levels}
TODO: Compare persons, families, and organizations as complex adaptive systems while preserving their differences.

\section{Brief Health Screening Logic}
TODO: Explain how screening logic travels across levels and where it must be adapted.

\section{Measurement Architecture}
TODO: Present the cross-level architecture that supports the OHA and the staged calibration design.

\section{Implications}
TODO: Discuss implications for instrument development, simulation, and validation planning.

\section{Conclusion}
TODO: State the theoretical contribution and its relationship to the dissertation.

\printbibliography

\end{document}
